\documentclass[12pt]{article}

\usepackage{fullpage}
\usepackage{graphicx}
\usepackage{amssymb}
\usepackage{amsmath}
\usepackage[none]{hyphenat}
\usepackage{parskip}
\usepackage[spanish]{babel}
\usepackage[utf8]{inputenc}
\usepackage{hyperref}
\usepackage{fancyhdr}
\usepackage{tasks}
\usepackage{mdframed}
\usepackage{xcolor}
\usepackage{pgfplots}
\usepackage[makeroom]{cancel}
\usepackage{multicol}
\usepackage[shortlabels]{enumitem}
\usepackage{stackrel}
\usepackage{tkz-tab}
\usepackage{xpatch}
\xpatchcmd{\tkzTabLine}{$0$}{$\bullet$}{}{}

\setlength{\headheight}{10pt}
\setlength{\headsep}{10pt}
\pagestyle{fancy}
\rhead{\ayudantia \ - \alumno}
\tikzset{t style/.style={style=solid}}

\newcommand*{\mybox}[2]{\colorbox{#1!30}{\parbox{.98\linewidth}{#2}}}

\newenvironment{solucion}
{\begin{mdframed}[backgroundcolor=black!10]
		{\bf Solución:}\\
	}
	{
	\end{mdframed}
}

\newenvironment{alternativas}[1]
{\begin{multicols}{#1}
		\begin{enumerate}[a)]
		}
		{
		\end{enumerate}
	\end{multicols}
}

\newenvironment{preguntas}
{\begin{enumerate}\itemsep12pt
	}
	{
	\end{enumerate}
}

\newcommand{\ayudantia}{{\sc Ayudantía 13}}
\newcommand{\tituloayu}{Ortogonalidad y formas cuadráticas}
\newcommand{\fecha}{13 de junio de 2019}
\newcommand{\sigla}{MAT1203}
\newcommand{\nombre}{Álgebra Lineal}
\newcommand{\profesor}{Camilo Perez}
\newcommand{\ano}{2019}
\newcommand{\semestre}{1}
\newcommand{\mail}{mat1203@ifcastaneda.cl}
\newcommand{\alumno}{Ignacio Castañeda - \mail}

\newcommand{\ev}{\Big|}
\newcommand{\ra}{\rightarrow}
\newcommand{\lra}{\leftrightarrow}
\newcommand{\N}{\mathbb{N}}
\newcommand{\R}{\mathbb{R}}
\newcommand{\Exp}[1]{\mathcal{E}_{#1}}
\newcommand{\List}[1]{\mathcal{L}_{#1}}
\newcommand{\EN}{\Exp{\N}}
\newcommand{\LN}{\List{\N}}
\newcommand{\comment}[1]{}
\newcommand{\lb}{\\~\\}
\newcommand{\eop}{_{\square}}
\newcommand{\hsig}{\hat{\sigma}}
\newcommand{\widesim}[2][1.5]{
	\mathrel{\overset{#2}{\scalebox{#1}[1]{$\sim$}}}
}
\newcommand{\wsim}{\widesim{}}
\newcommand{\lh}{\stackrel{L'H}{=}}

\begin{document}
\thispagestyle{empty}

\begin{minipage}{2cm}
	\includegraphics[width=2cm]{../../../../img/logo.pdf}
	\vspace{0.5cm}
\end{minipage}
\begin{minipage}{\linewidth}
	\begin{tabular}{lrl}
		{\scriptsize\sc Pontificia Universidad Catolica de Chile} & \hspace*{0.7in}Curso: &
		\sigla  - \nombre\\
		{\sc Facultad de Matemáticas}&
		Profesor: & \profesor \\
		{\sc Semestre \ano-\semestre} & Ayudante: & {Ignacio Castañeda}\\
		& {Mail:} & \texttt{\mail}
	\end{tabular}
\end{minipage}

\vspace{-10mm}
\begin{center}
	{\LARGE\bf \ayudantia}\\
	\vspace{0.1cm}
	{\tituloayu}\\
	\vspace{0.1cm}
	\fecha\\
	\vspace{0.4cm}
\end{center}

\begin{preguntas}
\item Sea 
	$U = Gen\left\{
	\begin{bmatrix}1\\0\\1\\1\end{bmatrix},
	\begin{bmatrix}0\\1\\2\\1\end{bmatrix},
	\begin{bmatrix}1\\1\\3\\2\end{bmatrix}
	\right\}$.
\begin{enumerate}[a)]
\item Encuentre una base ortonormal de $U$.
\item Encuentre la distancia de $b = 
		\begin{bmatrix}1\\1\\1\\1\end{bmatrix}$ a $U$.
\end{enumerate}
\begin{solucion}

\begin{enumerate}[a)]
\item Encuentre una base ortonormal de $U$.\\
			\\
			En primer lugar, debemos buscar una base común y corriente de $U$. Para esto, tomemos los vectores L.I. del conjunto que lo genera. Notemos que
			$$\begin{bmatrix}
			1 & 0 & 1 \\
			0 & 1 & 1 \\
			1 & 2 & 3 \\
			1 & 1 & 2
			\end{bmatrix} \sim
			\begin{bmatrix}
			1 & 0 & 1 \\
			0 & 1 & 1 \\
			0 & 0 & 0 \\
			0 & 0 & 0
			\end{bmatrix}$$
			Por lo que concluimos que los dos vectores del generado son L.I. Entonces, una base de $U$ corresponde a
			$$B = \left\{ \begin{pmatrix}
			1 \\ 0 \\ 1 \\ 1
			\end{pmatrix},\begin{pmatrix}
			0 \\ 1 \\ 2 \\ 1
			\end{pmatrix} \right\}$$
			Recordemos que Gramm-Schmidt forma una base ortogonal a partir de una base cualquiera y funciona de la siguiente forma:
			
			Sea $B = \{v_1, v_2, \dots \}$ una base cualquiera de un espacio vectorial,
			
			$$\begin{array}{rcl}
			u_1 & = & v_1 \\\\
			u_2 & = & v_2 - \dfrac{v_2u_1}{u_1u_1}u_1 \\\\
			u_3 & = & v_3 - \dfrac{v_3u_1}{u_1u_1}u_1 - \dfrac{v_3u_2}{u_2u_2}u_2\\
			& \vdots & 			
			\end{array}$$
			Luego, $B^{\perp} = \{u_1, u_2, \dots \}$ es una base ortogonal de ese espacio vectorial.
			
			Hagamos esto con nuestra base:
			
			En primer lugar,
			$$u_1 = v_1 = \begin{pmatrix}
			1 \\ 0 \\ 1 \\ 1
			\end{pmatrix}$$
			
			Luego,
			$$u_2 = v_2 - \dfrac{v_2u_1}{u_1u_1}u_1
			= \begin{pmatrix}
			0 \\ 1 \\ 2 \\ 1
			\end{pmatrix} - \dfrac{\begin{pmatrix}
				0 \\ 1 \\ 2 \\ 1
				\end{pmatrix}\begin{pmatrix}
				1 \\ 0 \\ 1 \\ 1
				\end{pmatrix}}{\begin{pmatrix}
				1 \\ 0 \\ 1 \\ 1
				\end{pmatrix}\begin{pmatrix}
				1 \\ 0 \\ 1 \\ 1
				\end{pmatrix}}\begin{pmatrix}
			1 \\ 0 \\ 1 \\ 1
			\end{pmatrix} 
			= \begin{pmatrix}
			0 \\ 1 \\ 2 \\ 1
			\end{pmatrix} - \dfrac{3}{3}\begin{pmatrix}
			1 \\ 0 \\ 1 \\ 1
			\end{pmatrix} 
			= \begin{pmatrix}
			-1 \\ 1 \\ 1 \\ 0
			\end{pmatrix} $$
			De esta forma,
			$$B^{\perp} = \left\{\begin{pmatrix}
			1 \\ 0 \\ 1 \\ 1
			\end{pmatrix}, \begin{pmatrix}
			-1 \\ 1 \\ 1 \\ 0
			\end{pmatrix}\right\}$$
			Sin embargo, necesitamos una base ortonormal. Para esto, basta con dividir todos los vectores de la base por su norma, obteniendo
			$$\hat{B}^{\perp} = \left\{\begin{pmatrix}
			1/\ \sqrt[]{3} \\ 0 \\ 1/\ \sqrt[]{3} \\ 1/\ \sqrt[]{3}
			\end{pmatrix}, \begin{pmatrix}
			-1/\ \sqrt[]{3} \\ 1/\ \sqrt[]{3} \\ 1/\ \sqrt[]{3} \\ 0
			\end{pmatrix}\right\}$$
\item Encuentre la distancia de $b = 
			\begin{bmatrix}1\\1\\1\\1\end{bmatrix}$ a $U$.
			
			Sea $A$ una matriz con los vectores de una base de $U$ en sus columnas y $x$ un vector coordenada de un elemento cualquier perteneciente a $U$, entonces la distancia entre ese elemento y $b$ corresponde a $||Ax-b||$.
			
			Luego, la distancia entre $b$ y $U$ corresponde a la distancia más pequeña entre $b$ y un elemento de $U$, es decir,
			$$min||Ax-b||$$
			Resolver este problema es equivalente a resolver el sistema
			$$A^TAx = A^Tb$$
			Utilicemos
			$$A = \begin{bmatrix}
			1 & 0 \\
			0 & 1 \\
			1 & 2 \\
			1 & 1
			\end{bmatrix}$$
			Entonces,
			$$A^TA = \begin{bmatrix}
			3 & 3 \\
			3 & 6
			\end{bmatrix}, \quad A^Tb = \begin{pmatrix}
			3 \\ 4
			\end{pmatrix}$$
			Por lo que tenemos que resolver el sistema
			$$\begin{bmatrix}
			3 & 3 \\
			3 & 6
			\end{bmatrix}x = \begin{pmatrix}
			3 \\ 4
			\end{pmatrix} \ra x = \begin{pmatrix}
			\frac{2}{3}\\ \frac{1}{3}
			\end{pmatrix}$$
			Finalmente, la distancia corresponde a
			$$||Ax-b|| = \left|\left|\begin{bmatrix}
			1 & 0 \\
			0 & 1 \\
			1 & 2 \\
			1 & 1
			\end{bmatrix}\begin{pmatrix}
			\frac{2}{3}\\ \frac{1}{3}
			\end{pmatrix} - \begin{pmatrix}1\\1\\1\\1\end{pmatrix}\right|\right|
			= \left|\left|\begin{pmatrix}
			-\frac{1}{3}\\
			-\frac{2}{3}\\
			\frac{1}{3}\\
			0
			\end{pmatrix}\right|\right| = \sqrt[]{\dfrac{2}{3}}$$
\end{enumerate}
\end{solucion}
\item Un cierto experimento genera los datos $(1,3)$, $(2,5)$ y $(3,4)$. Describa el modelo que da un ajuste de mínimos cuadrados de esos puntos mediante una recta de la forma $y = \beta_0 +  \beta_1x$
\begin{solucion}
Los valores del experimento se pueden ver en las siguiente tabla:\\
		\\
		\begin{center}
		\begin{tabular}{|l|l|}
			\hline
			\textbf{y} & \textbf{x} \\ \hline
			3          & 1          \\ \hline
			5          & 2          \\ \hline
			4          & 3          \\ \hline
		\end{tabular}
		\end{center}
		Estamos buscando $\beta_0, \beta_1$ tal que
		$$y = \beta_0 +  \beta_1x$$
		sea un ajuste por mínimos cuadrados. Esto lo podemos escribir como
		$$Y = X\beta$$
		donde
		$$Y = \begin{pmatrix}
		3 \\ 5 \\ 4
		\end{pmatrix}, \quad 
		X = \begin{bmatrix}
		1 & 1 \\
		1 & 2 \\
		1 & 3
		\end{bmatrix}, \quad
		\beta = \begin{pmatrix}
		\beta_0 \\ \beta_1
		\end{pmatrix}$$
		Para encontrar $\beta$, debemos resolver el sistema
		$$X^TX\beta = X^TY \ra \beta = (X^TX)^{-1}X^TY$$
		$$(X^TX)^{-1} = \left(\begin{bmatrix}
		1 & 1 & 1\\
		1 & 2 & 3
		\end{bmatrix}
		\begin{bmatrix}
		1 & 1 \\
		1 & 2 \\
		1 & 3
		\end{bmatrix}\right)^{-1}
		 = 
		 \begin{bmatrix}
		 3 & 6 \\
		 6 & 14
		 \end{bmatrix}^{-1} = \begin{bmatrix}
		 \frac{7}{3} & -1 \\
		 -1 & \frac{1}{2}
		 \end{bmatrix}$$
		 $$ X^TY = \begin{bmatrix}
		 1 & 1 & 1\\
		 1 & 2 & 3
		 \end{bmatrix} \begin{pmatrix}
		 3 \\ 5 \\ 4
		 \end{pmatrix} = 
		 \begin{pmatrix}
		 12 \\ 25
		 \end{pmatrix}$$
		 Luego,
		 $$\beta = (X^TX)^{-1}X^TY
		 = \begin{bmatrix}
		 \frac{7}{3} & -1 \\
		 -1 & \frac{1}{2}
		 \end{bmatrix}\begin{pmatrix}
		 12 \\ 25
		 \end{pmatrix}
		 = \begin{pmatrix}
		 3 \\ \frac{1}{2}
		 \end{pmatrix}
		 $$
		 Finalmente, el ajuste por mínimos cuadrados es
		 $$y = 3 + \dfrac{1}{2}x$$
\end{solucion}
\item Sea $A$ la matriz simétrica que representa a la forma cuadrática
	$$Q(x) = 9x_1^2 + 7x_2^2 + 11x_3^2 - 8x_1x_2 + 8x_1x_3$$
\begin{enumerate}[a)]
\item Encuentre una matriz ortogonal $P$ tal que el cambio de variable $x = Py$ transforma la forma $x^TAx$ en una forma cuadrática sin productos cruzados.
\item Escriba la forma cuadrática en las nuevas variables y clasifíquela.
\end{enumerate}
\begin{solucion}

\begin{enumerate}[a)]
\item Encuentre una matriz ortogonal $P$ tal que el cambio de variable $x = Py$ transforma la forma $x^TAx$ en una forma cuadrática sin productos cruzados.\\
\\
En primer lugar, podemos ver que
$$A = \begin{bmatrix}
9 & -4 & 4\\
-4 & 7 & 0\\
4 & 0 & 11
\end{bmatrix}$$
Notemos que para que una forma cuadrática no tenga productos cruzados, su matriz asociada debe ser diagonal.\\

Luego, haciendo el cambio de variable,
$$x^TAx = (Py)^TA(Py) = y^TP^TAPy$$
Por lo tanto, para cumplir lo que nos piden, debe ocurrir que $P^TAP = D$ con $D$ una matriz diagonal.\\

Notemos que
$$P^TAP = D \ra A = PDP^T$$
Por lo que podemos encontrar $P$ diagonalizando ortogonalmente $A$.\\

Buscando los valores y vectores propios de $A$ obtendremos
$$\lambda_1 = 3 \ra v_1 = \begin{pmatrix}
2 \\ 2 \\ -1
\end{pmatrix},
\quad
\lambda_2 = 9 \ra v_2 = \begin{pmatrix}
1 \\ -2 \\ -2
\end{pmatrix}, 
\quad
\lambda_3 = 15 \ra v_3 = \begin{pmatrix}
2 \\ -1 \\ 2
\end{pmatrix}$$
Como queremos diagonalizar ortogonalmente, normalizamos cada vector propio, con lo que tenemos
$$v_1 = \begin{pmatrix}
\frac{2}{3} \\ \frac{2}{3} \\ -\frac{1}{3}
\end{pmatrix},
\quad
v_2 = \begin{pmatrix}
\frac{1}{3} \\ -\frac{2}{3} \\ -\frac{2}{3}
\end{pmatrix},
\quad
v_3 = \begin{pmatrix}
\frac{2}{3} \\ -\frac{1}{3} \\ \frac{2}{3}
\end{pmatrix} $$
Finalmente,
$$P = \begin{bmatrix}
\frac{2}{3} & \frac{1}{3} & \frac{2}{3} \\
\frac{2}{3} & -\frac{2}{3} & -\frac{2}{3} \\
-\frac{1}{3} & -\frac{1}{3} & \frac{2}{3}
\end{bmatrix}$$
\item Escriba la forma cuadrática en las nuevas variables y clasifíquela.\\
\\
De antes, tenemos que
$$x^TAx = (Py)^TA(Py) = y^TP^TAPy = y^TDy$$
Donde $D$ también era la matriz diagonal de la diagonalización ortogonal de $A$, por lo que
$$D = \begin{bmatrix}
3 & 0 & 0 \\
0 & 9 & 0 \\
0 & 0 & 15
\end{bmatrix}$$
Por lo tanto, como todos los elementos de la diagonal son positivos, la forma cuadrática es positiva definida.
\end{enumerate}
\end{solucion}
\end{preguntas}
\end{document}